\medskip
\noindent\emph{证明.}
由于 $E(X)\in\mathbb Z_p[[X]]$, (4.2.1) 由定义得出. 对于 (4.2.2), 有

\begin{equation}\tag{4.2.4}
\log(\Theta(1))=0
\qquad\text{由 (4.1.5) 中对 $\pi$ 的选取,}
\end{equation}

及

\begin{equation}\tag{4.2.5}
\operatorname{ord}(\Theta(1)-1)=\frac1{p-1}
\qquad\text{由 (4.2.1).}
\end{equation}

由 (4.2.4), 存在整数 $n\geq0$ 使 $\Theta(1)^{p^n}=1$
(这是对数之核的刻画), 而 (4.2.5) 蕴涵 $n=1$.

至于 (4.2.3.1), 只要知道 $\sum_{j=0}^{a-1}t^{p^j}\in\mathbb Z_p$,
取对数后立即计算即可得到. 而 $t^{p^j}$, $0\leq j\leq a-1$,
恰为 $t\in W(\mathbb F_{p^a})$ 在 $\mathbb Z_p$ 上的共轭元; 故

\[
\sum_{j=0}^{a-1}t^{p^j}
=\operatorname{Tr}_{W(\mathbb F_{p^a})/\mathbb Z_p}(t)\in\mathbb Z_p.
\]

在 (4.2.3.1) 两端令 $X=1$, 得到以下结果.

\subsubsection*{推论 4.3.}
设 $t\in\Omega$ 满足 $t^{p^a}=t$, 以 $\bar t\in\mathbb F_{p^a}$ 表示其剩余类. 则

\begin{equation}\tag{4.3.1}
\prod_{j=0}^{a-1}\Theta(t^{p^j})
=\xi_p^{\operatorname{tr}_{\mathbb F_{p^a}/\mathbb F_p}(\bar t)}.
\end{equation}

\medskip
\noindent\emph{底本注.}
印本在引理~4.2 中有三处可相互看出的记号失误:
在 (4.2.1) 的系数展开中写作 $A_1=\pi$;
将 Teichmüller 代表元置于 $W(\mathbb F_{q^a})$,
但 $t^{p^a}=t$ 及推论~4.3 均使用 $\mathbb F_{p^a}$;
在 (4.2.3.1) 的指数中印作 $\sum t^j$,
但乘积及随后的证明均使用 $\sum t^{p^j}$.
照录法文层保留这些形式; 校订法文层和英文层采用 $B_1$、$\mathbb F_{p^a}$ 及 $t^{p^j}$.
